\section{Conclusion}

The aim of this project was to implement the AMOEBA polarizable force field for high pressure 
water ice and assess its suitability for these types of classical MD simulations.
Polarizable force fields have not been widely used because of the amount of computational 
resources required to iteratively calculate induced dipoles.
Computational resources, however, have improved in recent years, making more complex 
force fields practical.
This approach is still generally favoured over \textit{ab initio} simulations for 
large systems which still face significant computational difficulties.

\par
High-pressure ice is well suited for benchmarking the accuracy of a novel potential 
due to the complex interactions present, as governed by the Bernal-Fowler ice rules.
The order/disorder Ice--VII/--VIII phase transition is important in the field of high-pressure
ice, occupying a significant portion of the phase diagram abouve 2 GPa.
Furthermore, the unique properties of Ice--VII allow a high-pressure transformation into
a `plastic' ice where the hydrogen atoms are able to rotate about the lattice points.

\par
AMOEBA was able to reproduce the VII/VIII phase transition at 3 GPa, albeit with
significant hysteresis.
Inspection of the hydrogen mean-squared displacement of the system just below
the simulated transition temperature revealed that some molecules had
begun to ondergo a reorientation.
This suggests that, given more time simulation time, the observed hysteresis may be
reduced.
The complexity of this potential, however, may prohibit such simulations due to the 
sheer quantity of computational power required.
At higher pressures and higher temperatures, AMOEBA Ice--VII displayed hydrogen rotation.
Whether this rotation can be classified as the `plastic' phase, or simply a relaxation of 
the Ice--VII structure from some metastable state is unclear.
In either case, the AMOEBA ice was able to melt.

\par
In future research, it may be interesting to see whether AMOEBA Ice--VII would be able to 
transition into ordered Ice--VIII which is known to involve a complicated kinetic pathway.
Moreover, detailed calculations on high-pressure, high-temperature runs might reveal more
information about plastic Ice--VII.
